\documentclass[11pt,a4paper]{article}
\usepackage[margin=0.90in]{geometry}
\usepackage[T1]{fontenc}
\usepackage{lmodern,microtype}
\usepackage{amsmath,amssymb,mathtools,amsthm}
\usepackage{booktabs,tabularx,array}
\usepackage{algorithm,algpseudocode}
\usepackage{enumitem}
\usepackage{xcolor}
\usepackage[hidelinks]{hyperref}
\usepackage{xurl}
\usepackage{bm}
\usepackage{listings}
\usepackage{fancyhdr}
\usepackage{titlesec}
\definecolor{navy}{HTML}{17324F}
\definecolor{muted}{HTML}{546474}
\definecolor{light}{HTML}{F1F5F8}
\hypersetup{colorlinks=true,linkcolor=navy,urlcolor=navy,citecolor=navy}
\setlist[itemize]{leftmargin=*,itemsep=2pt,topsep=4pt}
\setlist[enumerate]{leftmargin=*,itemsep=2pt,topsep=4pt}
\setlength{\parskip}{3pt}
\setlength{\parindent}{0pt}
\pagestyle{fancy}\fancyhf{}\fancyhead[L]{\small\color{muted} Adaptive BlockFlow --- technical specification}\fancyhead[R]{\small\color{muted} Working draft, October 2026}\fancyfoot[C]{\thepage}
\titleformat{\section}{\large\bfseries\color{navy}}{\thesection}{0.7em}{}
\titleformat{\subsection}{\normalsize\bfseries\color{navy}}{\thesubsection}{0.7em}{}
\newcommand{\C}{\mathcal C}
\newcommand{\A}{\mathcal A}
\newcommand{\Sset}{\mathcal S}
\newcommand{\V}{\mathcal V}
\newcommand{\E}{\mathbb E}
\newcommand{\N}{\mathcal N}
\newcommand{\Cat}{\operatorname{Cat}}
\newcommand{\softmax}{\operatorname{softmax}}
\newcommand{\stopgrad}{\operatorname{sg}}
\newcommand{\NFE}{\operatorname{NFE}}
\newcommand{\CE}{\operatorname{CE}}
\newcommand{\mask}{\operatorname{mask}}
\newcommand{\KL}{\operatorname{KL}}
\lstdefinestyle{compact}{basicstyle=\ttfamily\small,backgroundcolor=\color{light},frame=single,rulecolor=\color{muted!20},breaklines=true,breakatwhitespace=true,columns=fullflexible,showstringspaces=false}
\lstset{style=compact}
\title{\vspace{-1.0cm}\textbf{Adaptive BlockFlow}\\[3pt]
\large All-Remaining-Token Categorical Flow Matching\\with Reinforcement-Learned Token Revelation\\[9pt]
\normalsize Technical Specification and Implementation Contract}
\author{Research working draft --- for review and implementation}
\date{October 2026 \quad | \quad Version 4: Gaussian-to-categorical VFM}
\begin{document}
\maketitle
\vspace{-0.7cm}
\begin{abstract}
We propose \emph{Adaptive BlockFlow}, a two-stage conditional translation model. Unlike conventional block diffusion, the generator predicts \emph{every as-yet-uncommitted target position} during each generation round. Stage~1 commits only a fixed number $b$ of these candidate tokens; Stage~2 learns an adaptive subset-selection policy via group-relative policy optimization (GRPO). The generative process adopts the \emph{Gaussian-to-categorical} formulation of Categorical Flow Maps: data tokens are represented as one-hot vocabulary vectors, noisy trajectories are \emph{continuous}, and a transformer predicts a categorical endpoint distribution supervised by variational flow matching (VFM). We use Set Diffusion's set-causal attention and incremental-caching ideas as an architectural starting point, and the official Categorical Flow Maps text implementation as a reference for the Gaussian-input projection, VFM objective, and solver. This report describes a proposed algorithm, not a claim that the published repositories already provide its combined implementation. No flow-map self-distillation, one-step distillation, or few-step distillation objective is included.
\end{abstract}
\vspace{0.2cm}
\tableofcontents
\newpage

\section{Non-negotiable algorithmic contract}
\textbf{Core rule:} on revelation round $k$, run conditional flow generation on the \emph{entire remaining target set} $\A_k$. The generator outputs candidate tokens for all $i\in\A_k$. Commit only $\Sset_k\subseteq\A_k$ and discard the other candidates. Previously committed tokens never change.
\begin{center}
\begin{tabularx}{\linewidth}{@{}p{0.26\linewidth}X@{}}\toprule
Element & Required behavior \\ \midrule
Task & WMT14 German $\rightarrow$ English; source sentence fixed and always visible. \\
Noise & Independent Gaussian noise per \emph{uncommitted} target position, $z_0^i\sim\N(0,I_V)$; not a $[\mathrm{MASK}]$ token or a uniformly sampled vocabulary ID. \\
Prediction & Categorical distributions over vocabulary items at \emph{all} $i\in\A_k$. \\
Stage 1 & Train with conditional VFM over all $\A_k$; inference commits $\min(b,|\A_k|)$ positions in a predetermined order. \\
Stage 2 & Freeze generator initially; a learned policy selects both subset cardinality and token positions with GRPO. \\
Attention & All active positions may interact; committed context must not depend on unrevealed noisy positions. \\
Caching & Functional correctness first, uncached; cache only invariant clean-context states after parity tests. \\
Scope exclusion & No flow-map self-distillation or one-/few-step distillation training. \\ \bottomrule
\end{tabularx}
\end{center}
\textbf{Terminology:} the training data and endpoint predictions are categorical, but the intermediate flow states $z_t$ are \emph{continuous} in $\mathbb R^V$. Thus our chosen model is a \emph{continuous Gaussian-to-categorical variational flow} for discrete text, not a discrete-state continuous-time Markov chain. This follows the underlying VFM construction of Roos et al.~\cite{roos2026}.

\section{Motivation and relation to prior work}
Standard Block Diffusion generates and denoises a current block, then moves to the next. Our Stage~1 deliberately differs: it solves a flow problem over \emph{all remaining positions} each round, even though it adds only $b$ tokens to context. Stage~2 learns the subset selection rule. Set Diffusion \cite{arriola2026} supplies relevant set-causal ordering and KV-cache ideas, but the existing implementation is masked discrete diffusion, not our continuous Gaussian-to-categorical generator. Categorical Flow Maps \cite{roos2026} supplies the flow path, endpoint prediction, and VFM loss, but does not supply the proposed revelation-controller training and WMT14 conditional architecture.

\paragraph{Software references (not drop-in compatibility).}
\begin{itemize}
\item \textbf{Set Diffusion:} \url{https://github.com/kuleshov-group/setdlms}. Reuse/refine its backbone, set-order masks, cache logic, and evaluation infrastructure after inspecting real APIs.
\item \textbf{Categorical Flow Maps (text implementation):} \url{https://github.com/olsdavis/semicat}. Inspect \texttt{semicat/models/semicat.py} and \texttt{semicat/models/textsemicat.py} for Gaussian priors, continuous-input embeddings, time conditioning, VFM, and basic sampling; disable the repository's self-distillation path for our initial experiments.
\item \textbf{Training from scratch:} no ELF checkpoint; no assumption that a Set Diffusion or Semicat pretrained checkpoint is compatible with the combined architecture.
\end{itemize}

\section{Notation and conditional generation semantics}
Given a German sentence $s$, let $y=(y_1,\ldots,y_L)$ be the target English token sequence including an EOS token, with token alphabet $\V$ of size $V$. Denote one-hot token vectors by $e(y_i)\in\{e_1,\ldots,e_V\}\subset\mathbb R^V$.
\begin{center}
\begin{tabularx}{\linewidth}{@{}lX@{}}\toprule
Symbol & Definition \\ \midrule
$\mathcal U=\{1,\dots,L_{\rm canvas}\}$ & Allocated target canvas; generated without reading reference length at test time. \\
$\C_k$ & Permanently committed target positions before round $k$. \\
$\A_k=\mathcal U\setminus\C_k$ & All remaining target positions on round $k$. \\
$\Sset_k\subseteq\A_k$ & Positions newly committed at the end of round $k$. \\
$z_t^{\A_k}$ & Continuous noisy states for all remaining positions at flow time $t$. \\
$p_\theta^i(\cdot\mid z_t^{\A_k},t,s,y_{\C_k})$ & Categorical endpoint distribution for position $i\in\A_k$. \\
$\widehat y_{\A_k}$ & Candidate token IDs sampled/decoded from final flow endpoints. \\
$\pi_\phi$ & RL policy over nonempty subsets of $\A_k$ (and therefore their sizes). \\
$T$ & Number of network evaluations/solver steps in one complete flow solve (subject to solver details). \\
\bottomrule
\end{tabularx}
\end{center}
For training, $\mathcal U$ contains only valid reference positions up to EOS (plus padding excluded from loss). During inference, $L_{\rm canvas}$ is chosen by a source-only length convention, never from the reference translation. We must explicitly reconcile variable-length output with EOS and canvas limits (Section~\ref{sec:length}).

\section{Gaussian-to-categorical variational flow matching}
\subsection{Noise and probability path}
For an active target position $i\in\A$ with clean one-hot endpoint $x_1^i=e(y_i)$, draw independent Gaussian noise $x_0^i\sim\N(0,I_V)$ and time $t\sim\mathrm{Uniform}(0,1)$. Define the \emph{continuous} interpolant
\begin{equation}
\boxed{z_t^i=(1-t)x_0^i+t e(y_i),\qquad i\in\A.}\label{eq:interpolant}
\end{equation}
The source and committed tokens remain clean; do \emph{not} interpolate them and do \emph{not} allow their representations to absorb information from the current noisy target positions. The model predicts endpoint distributions
\begin{equation}
\boxed{p_\theta^i=\softmax\!\left(f_\theta(z_t^{\A},t,s,y_{\C})_i\right)\in\Delta^{V-1},\quad i\in\A.}\label{eq:endpoint}
\end{equation}
This produces vocabulary logits through a standard vocabulary projection. It is not a separate, second text decoder network. The \emph{input} differs from ordinary discrete-token diffusion: uncommitted states are continuous $V$-dimensional vectors, requiring a dedicated stable continuous-input projection to transformer hidden width $d$ (informed by Appendix~G.1 of \cite{roos2026}). The \emph{clean source/committed} tokens can use normal token-embedding lookup.

\subsection{Conditional VFM objective}
The VFM endpoint loss over \emph{all remaining target positions} is
\begin{equation}
\boxed{
\mathcal L_{\mathrm{VFM}}(\theta)=
\E_{(s,y),\C,t,x_0}\left[-\frac{1}{|\A|}\sum_{i\in\A}
\log p_\theta^i\left(y_i\mid z_t^{\A},t,s,y_{\C}\right)\right],
\qquad \A=\mathcal U\setminus\C.}
\label{eq:loss}
\end{equation}
Padding, source tokens, and committed tokens are excluded from the target loss. For all valid $i\in\A$, the supervision is used, even if $i$ will \emph{not} be among the next $b$ committed positions. No self-distillation loss is needed to define or train this VFM model.

\subsection{Instantaneous-flow sampler (no distillation)}
The endpoint prediction induces the continuous velocity field (for $t<1$)
\begin{equation}
\boxed{v_\theta(z_t,t\mid s,y_{\C})
=\frac{p_\theta(z_t,t\mid s,y_{\C})-z_t}{1-t}.}\label{eq:velocity}
\end{equation}
With grid $0=t_0<\cdots<t_N=1$, an explicit Euler step for remaining positions is
\begin{equation}
z_{t_{n+1}}^{\A}=z_{t_n}^{\A}+
\frac{t_{n+1}-t_n}{1-t_n}
\left(p_\theta(z_{t_n}^{\A},t_n\mid s,y_\C)-z_{t_n}^{\A}\right).
\label{eq:euler}
\end{equation}
The last step of this schedule transports $z_{t_{N-1}}$ to the predicted simplex-valued endpoint. Sample candidate token IDs as $\widehat y_i\sim\Cat(z_1^i)$ (or use argmax as a deterministic baseline). Numerical details must be validated against the VFM-only solver, with stable denominators and explicit reporting of the actual number of network calls. Equation~\eqref{eq:euler} is an \emph{instantaneous velocity} sampler, not a learned interval flow-map or a distillation method.

\section{Stage 1: all-remaining-token generation with fixed commitment}
\subsection{Context distribution during training}
Let $b\ge1$ be the fixed commitment size. For left-to-right training, sample a reveal-round index $k$ and form
\begin{equation}
\C_k=\{1,\ldots,\min(kb,L)\},\qquad
\A_k=\{1,\ldots,L\}\setminus\C_k.
\end{equation}
On the resulting context, construct Gaussian-to-data paths for \emph{every} $i\in\A_k$ using Equation~\eqref{eq:interpolant}, and optimize Equation~\eqref{eq:loss}. The context is teacher-forced clean target tokens during initial supervised training. It must not receive information from the unrevealed reference tokens via attention or cached hidden states.


\begin{algorithm}[htbp]
\caption{Stage 1 training}\label{alg:train}
\begin{algorithmic}[1]
\Require The dataset $\mathcal{D}$, size $b$, denoiser $p_\theta$
\State Choose left-to-right ordering.
\State Sample $\bm{x} = [\bm{x}_1, \ldots, \bm{x}_{L}] \sim \mathcal{D}$
\State For each step $k$; obtain $\C_k = \{1,\ldots, kb\}$ and $\A_k= \{kb+1, L\}$
\State Sample $t\sim U(0,1)$ and $\bm{\epsilon}^i\sim\N(0,I_V)$ for \textbf{all} $i\in\A_k$
\State $x_i(t)\gets(1-t)\bm{\epsilon}^i+t (\bm{x}_i)$ for \textbf{all} $i\in\A_k$
\State Predict vocabulary distributions $p_\theta^i$ at all $i\in\A_k$
\State $\ell\gets-|\A_k|^{-1}\sum_{i\in\A_k}\log p_\theta^i(x_i\mid x_i(t), \bm{x}_{c})$
\State Update $\theta$ using $\nabla_\theta\ell$.
\end{algorithmic}
\end{algorithm}

\subsection{Stage 1 inference: fixed $b$ committed each round}
A \emph{complete flow solve} is performed for $\A_k$ on every round, from fresh Gaussian noise (initial baseline). Once candidate IDs for all $i\in\A_k$ are obtained, commit exactly the next $\min(b,|\A_k|)$ valid positions and discard candidates elsewhere. Later rounds predict all remaining positions again.
\begin{algorithm}[htbp]
\caption{Stage 1 inference: all remaining flow, fixed commitment}\label{alg:infer}
\begin{algorithmic}[1]
\Require commit size $b$, flow steps $T$
\State $\C\gets\varnothing$; committed tokens $\bm{x}_c\gets\varnothing$, $\A \gets \{1, \ldots,L \}$
\While{$\A\ne\varnothing$}
  \State Sample $\bm{x}_i(0)\sim\N(0,I_V)$ for \textbf{all} $i\in\A$
  \State Update flow on \textbf{all} $\A$ for $T$ evaluations
  \begin{align}
      \nonumber \bm{x}_{i}(t + \Delta t) = \bm{x}_i(t) + \Delta t \frac{p_{\theta}(\bm{x}_i \mid \bm{x}_i(t), t, \bm{x}_c) - \bm{x}_i(t)}{1 - t}
  \end{align}
  \State Decode candidate token IDs $ \hat{\bm{x}}_i$ for \textbf{all} $i\in\A$
  \State $\Sset\gets$$\min(b,|\A|)$ positions in $\A$ (fixed left to right)
  \State Commit $\bm{x}_i\gets\hat{\bm{x}}_i$ only for $i\in\Sset$
  \State $\C\gets\C\cup\Sset$; $\A\gets\A\setminus\Sset$
\EndWhile
\State \Return $\bm{x}(1)$
\end{algorithmic}
\end{algorithm}

\paragraph{Example ($L=12$, $b=3$).}
Round~0 flows over positions $1{:}12$, commits $1{:}3$; round~1 flows over $4{:}12$, commits $4{:}6$; round~2 flows over $7{:}12$, commits $7{:}9$; round~3 flows over $10{:}12$, commits $10{:}12$. The predictions at positions outside the committed subset are discarded after each round. This is \emph{not} the sampler that denoises only $1{:}3$, then only $4{:}6$, etc.

\section{Stage 2: GRPO for adaptive token revelation}
\subsection{Policy state and action}
Freeze $\theta$ initially. At reveal round $k$ the \emph{same} generator solves the flow over $\A_k$ and exposes candidate token IDs and endpoint probabilities. The policy observation $o_k$ includes:
\begin{equation}
o_k=(h(s),\;\C_k,\;\A_k,\;\widehat x_{\A_k},\;\text{entropy}_{\A_k},\;
p_{\theta}(\bm{x}_{i \in \A} \mid \bm{x}_{c})).
\end{equation}
Policy features from generator hidden states are detached. Reference targets are used only to calculate training-time reward. The policy samples a \emph{nonempty subset} $\Sset_k\subseteq\A_k$, and its size $m_k=|\Sset_k|$ is adaptive. A convenient implementable factorization is
\begin{align}
 m_k&\sim p_\phi(m\mid o_k),\quad 1\le m\le |\A_{k}|,\\
 (i_1,\ldots,i_{m_k})&\sim \prod_{r=1}^{m_k}p_\phi(i_r\mid o_k,m_k,i_{<r}),
\qquad i_r\in\A_k\setminus\{i_1,\ldots,i_{r-1}\}.\label{eq:act}
\end{align}
The action is the sampled \emph{ordered sequence} $(m_k,i_1,\ldots,i_{m_k})$ whose underlying subset is committed. Its exact log probability is
\begin{equation}
\boxed{\log\pi_\phi(a_k\mid o_k)=\log p_\phi(m_k\mid o_k)
+\sum_{r=1}^{m_k}\log p_\phi(i_r\mid o_k,m_k,i_{<r}).}\label{eq:logp}
\end{equation}
No duplicate selection is possible. If the policy is defined over \emph{unordered} sets instead, an ordering aggregation is necessary; do not claim Equation~\eqref{eq:logp} is the normalized log probability of an unordered subset. For policy gradients, storing the actual sampled ordered action is valid because that is the chosen action space.

\subsection{Rewards: quality, diversity, computational cost}
Let $\tau_g$ denote one complete rollout for a German source $s$, producing English output $\widehat y_g$. Use a group of $G\ge2$ independent rollouts for the \emph{same} source. Set
\begin{equation}
\boxed{R_g= Q_g+\lambda_d D_g-\lambda_c C_g,}\label{eq:reward}
\end{equation}
\textbf{For the WMT14 task}. Set $\lambda_{d} = 0$, donot consider the diversity reward.
\begin{itemize}
\item \textbf{Quality $Q_g$:}
    Measure translation accuracy using sentence-level chrF:
    \[
    Q_g = \frac{\mathrm{chrF}(\widehat{\bm{y}}_g,\bm{y})}{100},
    \]
    where $\bm{y}$ is the reference translation.

    \item \textbf{Computational cost $C_g$:}
    Penalize the number of function evaluations (NFE),
    i.e., the number of generator forward passes:
    \[
    C_g = \frac{\mathrm{NFE}(\tau_g)}{N_{\mathrm{ref}}}.
    \]
    If each revelation round requires $T$ model evaluations,
    then $\mathrm{NFE}(\tau_g)=K_gT$, where $K_g$ is the
    number of revelation rounds.
\end{itemize}


\subsection{Group-relative policy optimization}
For each source, obtain $G$ rollouts under a frozen old controller $\pi_{\rm old}$. For terminal rewards $R_1,\ldots,R_G$ compute
\begin{equation}
\widehat A_g=\frac{R_g-\bar R}{\operatorname{std}(R_1,\ldots,R_G)+\epsilon_A},\qquad
\bar R=\frac1G\sum_{h=1}^G R_h.\label{eq:adv}
\end{equation}
If all rewards are identical, set the advantage to zero rather than divide by an unstable tiny value. At each stochastic reveal decision $k$, store $o_{g,k}$, ordered $a_{g,k}$, and $\log\pi_{\rm old}(a_{g,k}\mid o_{g,k})$. Define
\begin{equation}
 r_{g,k}(\phi)=\exp\!\left[\log\pi_\phi(a_{g,k}\mid o_{g,k})-
 \log\pi_{\rm old}(a_{g,k}\mid o_{g,k})\right].
\end{equation}
A basic clipped GRPO objective \emph{to maximize} is
\begin{equation}
\boxed{J_{\rm GRPO}(\phi)=\E\left[
\frac1G\sum_{g=1}^G \frac{1}{K_g}\sum_{k=0}^{K_g-1}
\min\!\left(r_{g,k}(\phi)\widehat A_g,\,
\operatorname{clip}(r_{g,k}(\phi),1-\varepsilon,1+\varepsilon)\widehat A_g\right)
\right]-\beta\,\mathcal R_{\rm ref}(\phi).}\label{eq:grpo}
\end{equation}
$\mathcal R_{\rm ref}$ is an \emph{optional} controller-reference regularizer; leave it off initially unless implemented and verified. Forced last-round actions must be deterministic and excluded from the stochastic policy loss. This objective updates only $\phi$; the generator $\theta$ stays frozen for the initial Stage~2 experiment. Using a common terminal advantage for every decision is a simple episodic baseline, not a claim of optimal credit assignment.

\begin{algorithm}[htbp]
\caption{Stage 2 rollout and policy update (frozen generator)}\label{alg:rl}
\begin{algorithmic}[1]
\Require Source $s$, frozen generator $\theta$, controller $\pi_\phi$, $G$ rollouts
\For{$g=1,\ldots,G$}
  \State Initialize empty context $\C$, all remaining positions $\A$
  \While{$\A\ne\varnothing$ and budget permits}
    \State Sample Gaussian noise and solve flow on \textbf{all} $\A$
    \State Decode all $\widehat y_{\A}$ and construct reference-free observation $o$
    \State Draw nonempty ordered action $a\sim\pi_\phi(\cdot\mid o)$; extract $\Sset$
    \State Save $o,a,\log\pi_{\rm old}(a\mid o)$
    \State Commit only $\widehat y_{\Sset}$; update $\C,\A$ and source/context cache
  \EndWhile
  \State Record terminal translation $\widehat y_g$, actual NFE, latency and reward inputs
\EndFor
\State Compute group-relative advantages using Eq.~\eqref{eq:adv}
\State Update only $\phi$ by maximizing Eq.~\eqref{eq:grpo}
\end{algorithmic}
\end{algorithm}

\section{Attention, caching, and training--inference consistency}\label{sec:cache}
\subsection{Proposed attention contract}
The architecture must support \emph{active-to-active} interaction while keeping permanent context invariant. A sufficient uncached attention design is:
\begin{center}
\begin{tabularx}{\linewidth}{@{}p{0.26\linewidth}XXX@{}}\toprule
Query $\downarrow$, keys $\rightarrow$ & Source keys & Committed target keys & Active target keys \\ \midrule
Source queries & Allow & Disallow & Disallow \\
Committed queries & Allow & Allow only same/earlier reveal group & Disallow \\
Active queries & Allow & Allow all committed & Allow all active \\ \bottomrule
\end{tabularx}
\end{center}
For clean committed positions, causal structure is defined by \emph{revelation group}, not necessarily by absolute token index. In the first baseline, a full uncached forward pass is the ground truth implementation. Source/committed representations must be independent of the current flow time $t$; otherwise cached KV values become stale when $t$ changes. Time conditioning can be applied on an active-only branch (or another proved-equivalent mechanism).

\subsection{Safe incremental KV-cache construction}
\begin{enumerate}
\item Precompute clean source KV states using source-only attention.
\item Run each flow evaluation on all active positions with access to source and committed KV, and full active-to-active attention. \emph{Never} append changing noisy active states to permanent cache.
\item After choosing $\Sset_k$, convert its final candidates to hard token IDs; then perform a \emph{separate clean-context forward} for the newly committed group using no active-token keys. Append its invariant per-layer KV states.
\item Maintain absolute target positions and reveal-group metadata; do not assume that cache append order equals left-to-right target order.
\item Compare cached and fully recomputed logits/endpoint distributions under exactly the same sampled noise, target set, masks, times, and token IDs. Enable cache only after numerical parity tests pass.
\end{enumerate}
An implementation that allows a committed state to attend current noisy tokens, or injects changing flow-time embeddings globally into committed states, cannot safely reuse that cached state. Caching is therefore an architectural correctness constraint, not a claim that toggling \texttt{use\_cache=True} is sufficient.

\section{WMT14 setup and variable-length outputs}\label{sec:length}
Use the WMT14 German-to-English train/validation/test splits supported by the chosen dataset pipeline; fix the same tokenizer and scoring conventions across all baselines. Append target EOS during training. Keep source clean and visible; noise and VFM losses apply only to target $\A$. Report the exact split identity, preprocessing, and detokenization used.

\paragraph{Target length must be generated, not read from the reference.}
A first implementable approach is a \emph{source-only} length estimator or bucket predictor that proposes $L_{\rm canvas}$; alternatively choose a documented fixed maximum canvas and let a committed EOS terminate output when all earlier positions are committed. Compare these alternatives, record failed EOS and truncation. At arbitrary-order reveal, an EOS candidate at position $j$ cannot finish generation while any position $i<j$ remains uncommitted. Positions after the first valid committed EOS are ignored for decoding. The reference translation length, EOS index, and future target IDs must never enter inference or policy observations.

\paragraph{Important fairness note.}
Some legacy translation implementations use reference lengths at evaluation time; this would invalidate our claim of a practical generator. Log whether any baseline uses a supplied target length and do not compare wall-clock generation quality under unequal length information.

\section{Compute cost, memory, and experimental controls}
For a target canvas of length $L$ and fixed commitment size $b$, Stage~1 takes $K\approx\lceil L/b\rceil$ full-remaining flow solves. With $T$ network evaluations per solve, the approximate total denoiser count is $KT$; cache-context update calls and guidance may add evaluations. Stage~2 learns $K$ adaptively, while $T$ is initially fixed. The flow solves shrink in active length as positions commit, so NFE alone does not determine FLOPs or latency. Report all three.

Each continuous Gaussian active state is nominally $V$-dimensional. For batch size $B$, active length $M$, and float dtype bytes $q$, one dense $B\times M\times V$ buffer uses $qBMV$ bytes; multiple buffers, temporary projections, and logits can multiply peak memory. A robust design projects continuous active inputs to transformer width $d$, retains normal clean-token embedding lookup, and avoids materializing redundant $B\times M\times V$ arrays when possible. \emph{Do not} replace the specified Gaussian vocabulary-space prior with a low-dimensional Gaussian without documenting a distinct alternative model and loss.

\begin{center}
\begin{tabularx}{\linewidth}{@{}llX@{}}\toprule
ID & Method & Scientific question \\ \midrule
E0 & Full remaining, one solve & What is the no-revelation baseline at matched flow NFEs? \\
E1 & Stage 1 fixed $b\in\{4,8,16,32\}$ LTR & Quality/latency versus commitment size. \\
E2 & Stage 1 mixed fixed-size order & Does random-order context augmentation help Stage 2? \\
E3 & Random nonempty subset (budget matched) & Does any adaptation outperform an untrained selector? \\
E4 & Confidence/entropy-based selector & Is an RL controller better than a simple heuristic? \\
E5 & GRPO quality+cost & Main result, reward diversity disabled. \\
E6 & GRPO quality+cost+diversity & Optional reward ablation after E5 is stable. \\ \bottomrule
\end{tabularx}
\end{center}
Measure corpus BLEU and chrF, median/p95 translation latency, wall-clock throughput, number of reveal rounds, total generator NFEs, clean-cache rebuild calls, EOS failure rate, memory, and variation across seeds. Compare quality at \emph{matched NFE and matched latency budgets}; a controller that reduces $K$ but degrades translation quality may not improve the Pareto frontier.

\section{Implementation modules and execution milestones}
\begin{center}
\begin{tabularx}{\linewidth}{@{}p{0.39\linewidth}X@{}}\toprule
Proposed module (new, not a claim about repo APIs) & Contract \\ \midrule
\texttt{blockflow/data/translation.py} & WMT14 source-target tokenization, EOS, length, split. \\
\texttt{model/continuous.py} & Gaussian/continuous $V$-dimensional state projection and time embedding. \\
\texttt{blockflow/model/attention.py} & Source, reveal-group-committed, active set-causal masks. \\
\texttt{blockflow/flow/vfm.py} & Gaussian path, endpoint categorical CE, instantaneous velocity. \\
\texttt{blockflow/flow/sample.py} & Solve for \emph{all remaining} positions; return IDs, probabilities, NFE. \\
\texttt{blockflow/generate/fixed.py} & Fixed $b$ commit after complete all-remaining flow solve. \\
\texttt{blockflow/generate/cache.py} & Clean-token KV rebuild/append, uncached parity harness. \\
\texttt{blockflow/rl/policy.py} & Cardinality + ordered distinct-index sampling, valid log probability. \\
\texttt{blockflow/rl/grpo.py} & Group rollouts, task reward, clipped policy loss, frozen generator. \\
\texttt{blockflow/eval/translation.py} & chrF/BLEU, NFE, latency, EOS, output generation audits. \\ \bottomrule
\end{tabularx}
\end{center}
\textbf{M0 (inspect and preserve):} checkout both referenced repositories; record commit hashes, inspect true attention/cache model signatures, understand Semicat's VFM-only code path, and run a tiny baseline. Do not blindly merge their frameworks.\par
\textbf{M1 (VFM from scratch):} get source-conditioned Gaussian-to-categorical VFM training on a very small WMT14 subset. Verify finite loss and proper output simplex. No RL.\par
\textbf{M2 (fixed reveal, no cache):} implement all-remaining flow sampling, commit fixed $b$, ensure EOS correctness, prevent reference leakage. Establish nontrivial translation outputs.\par
\textbf{M3 (cache and heuristics):} add clean-context cache with cached/uncached parity, then random and confidence revelation baselines.\par
\textbf{M4 (GRPO):} freeze VFM generator, implement ordered subset policy, group rollouts, clipped GRPO, quality and cost reward; ablate diversity only afterward.\par
\textbf{M5 (full WMT14):} run controlled benchmarks and save configs, seeds, checkpoints, CPU/GPU/hardware characteristics, and evaluation artifacts.

\section{Mandatory correctness tests and failure modes}
\begin{enumerate}
\item \textbf{Noise path:} $z_0\sim\N(0,I_V)$, and $z_1=e(y)$ for every valid target position; no masked corruption, discrete-uniform vocabulary prior, or quiet change of model family.
\item \textbf{All-remaining objective:} nonzero supervised loss on all $i\in\A$, including those outside the next fixed $b$ commitment set; zero loss on committed/source/pad positions.
\item \textbf{Reference leakage:} replacing target IDs at unrevealed positions without changing their explicit interpolated noisy inputs cannot alter encoded clean source/committed context; no unobserved targets in policy observations.
\item \textbf{Noisy-state isolation:} clean context is independent of Gaussian seed and flow time, enabling safe cache reuse; full attention rule tested for noncontiguous selections.
\item \textbf{Commit immutability:} previously committed IDs never change, even when later flow solves predict different candidates elsewhere.
\item \textbf{Sampler validity:} predicted endpoints sum to one and remain nonnegative; final categorical sampling valid; fixed random seeds reproduce uncached generations.
\item \textbf{Fixed-$b$ semantics:} stage 1 always commits exactly $\min(b,|\A|)$ when no valid EOS shortcut applies; model still samples every remaining position at every round.
\item \textbf{KV parity:} cached/uncached outputs numerically agree at each time step for an identical fully specified generation trajectory and positional IDs.
\item \textbf{Policy log probability:} sampled ordered actions are nonempty, contain no duplicate/committed indices, and have correct normalized probability on tiny enumerable examples.
\item \textbf{GRPO:} gradient affects controller $\phi$ only; frozen $\theta$ stays unchanged; zero-variance groups and forced final actions handled correctly.
\item \textbf{EOS and length:} no reference-based target length at test time; termination well-defined for noncontiguous commitment.
\item \textbf{Compute accounting:} actual generator evaluations, context rebuilds and elapsed time measured rather than estimated or silently omitted.
\end{enumerate}

\section{Instructions for an implementation agent}
The source-of-truth contract is Sections~1, 4, 5, and~6. Begin by producing a short code-level map of \texttt{setdlms} and \texttt{semicat} showing exactly which functions can be reused. \emph{Inspect code and signatures before patching.} If a useful API is absent, say so and introduce a tested adapter rather than inventing its existence. Implement M0--M2 and execute smoke tests before any RL changes. Keep the original baseline runnable. Flag scientific decisions (choice of target canvas, joint versus per-token noise, absolute positions, fixed context mixture, duration $T$, and reward weights) explicitly in configuration and logs. Do not fabricate experiment metrics.

\section{Research claims and limitations}
The intended contribution is the \emph{generation/commitment decoupling} and a \emph{learned token-revelation policy} optimizing translation quality against inference cost. We do not claim an exact autoregressive likelihood interpolation, or that VFM automatically gives good one-step translation. Training on clean teacher-forced contexts can differ from inference on committed model predictions; arbitrary-order contexts and group-based reward diversity add further mismatch. All proposed benefits require controlled experiments. The named repositories are architectural/algorithmic references, not evidence that this new end-to-end system is already implemented.

\clearpage
\begin{thebibliography}{9}
\bibitem{roos2026} D. Roos et al., ``Categorical Flow Maps,'' \emph{International Conference on Machine Learning}, 2026. \url{https://proceedings.mlr.press/v306/roos26a.html}. Official text implementation: \url{https://github.com/olsdavis/semicat}.
\bibitem{arriola2026} M. Arriola and V. Kuleshov, ``Set Diffusion: Interpolating Token Orderings between Autoregression and Diffusion for Fast and Flexible Decoding,'' \emph{International Conference on Machine Learning}, 2026. \url{https://proceedings.mlr.press/v306/arriola26a.html}. Code: \url{https://github.com/kuleshov-group/setdlms}.
\bibitem{eijkelboom2024} F. Eijkelboom et al., ``Variational Flow Matching for Graph Generation,'' \emph{NeurIPS}, 2024.
\bibitem{gat2024} I. Gat et al., ``Discrete Flow Matching,'' \emph{NeurIPS}, 2024. (Different discrete-state formulation; contextual reference only.)
\bibitem{deepseekmath2024} Z. Shao et al., ``DeepSeekMath: Pushing the Limits of Mathematical Reasoning in Open Language Models,'' 2024. (Source of GRPO-style group-relative optimization.)
\end{thebibliography}
\end{document}